% Part: computability
% Chapter: computability-theory
% Section: universal-part-function

\documentclass[../../../include/open-logic-section]{subfiles}

\begin{document}

\olfileid{cmp}{thy}{uni}
\olsection{सार्वत्रिक आंशिक संगणनक्षम फलन}

\begin{thm}
\ollabel{thm:univ-comp}
एक सार्वत्रिक आंशिक संगणनक्षम फलन~$\fn{Un}(e,x)$ असते. दुसऱ्या
शब्दांत, पुढील गुणधर्म असलेले फलन $\fn{Un}(e,x)$ असते:
\begin{enumerate}
\item $\fn{Un}(e,x)$ आंशिक संगणनक्षम आहे.
\item $f(x)$ हे कोणतेही आंशिक संगणनक्षम फलन असेल, तर अशी नैसर्गिक
संख्या $e$ असते की प्रत्येक~$x$ साठी $f(x) \simeq \fn{Un}(e,x)$.
\end{enumerate}
\end{thm}

\begin{proof}
क्लिनी यांच्या प्रमाण रूप प्रमेयातील (\olref[nfm]{thm:normal-form})
$U$ आणि $T$ घेऊन $\fn{Un}(e,x) \simeq U(\umin{s}{T(e,x,s)})$
अशी व्याख्या करा.
\end{proof}

\begin{explain}
आंशिक संगणनक्षम फलनांचे प्रभावी प्रगणन आहे, हे अचूकपणे सांगण्याचा
हा मार्ग आहे. कल्पना अशी की $f_e(x) = \fn{Un}(e,x)$ ने परिभाषित
केलेल्या फलनाला $f_e$ म्हटले, तर $f_0$, $f_1$, $f_2$, \dots या
क्रमिकेत सर्व आंशिक संगणनक्षम फलने येतात; शिवाय $f_e(x)$ चे
संगणन $e$ आणि~$x$ या दोहोंवर ``एकसमानपणे'' करता येते. सोपेपणासाठी
आपण एकस्थानी फलनांसाठी सार्वत्रिक असलेले द्विस्थानी फलन वापरत आहोत.
परंतु संख्यांच्या क्रमिकांचे सांकेतिकरण करून हे अधिक आदानांसाठी सहज
सामान्य करता येते. उदाहरणार्थ, $f(x,y,z)$ हे $3$-स्थानी आंशिक
पुनरावर्ती फलन असेल, तर $g(x) \simeq f((x)_0, (x)_1, (x)_2)$
हे एकस्थानी आंशिक पुनरावर्ती फलन असते.
%READERNOTE{OLCMP-014 — गोठवलेल्या इंग्रजीत शेवटचे g हे केवळ ``recursive'' म्हटले आहे. f आंशिक पुनरावर्ती असल्याने g सुद्धा आंशिक पुनरावर्तीच हमखास असते; ते सर्वत्र परिभाषित असणे आवश्यक नाही. मराठीत हा गुणधर्म स्पष्ट केला आहे.}
\end{explain}

\end{document}
